\begin{comment}
\begin{tikzpicture}[
        x=1.15cm,
        y=1.05cm,
        >=Stealth,
        every node/.style={font=\small}
    ]

    % =========================================================
    % EIXOS
    % =========================================================

    \draw[->, thick]
        (0,0) -- (8.3,0)
        node[right] {Desvio-padrão $(\sigma_p)$};

    \draw[->, thick]
        (0,0) -- (0,6.5)
        node[above] {$E(R_p)$};

    \node[
        left,
        align=center
    ] at (-0.15,5.6)
        {Retorno\\esperado};


    % =========================================================
    % CONJUNTO DE CARTEIRAS POSSÍVEIS
    % =========================================================

    \path[
        fill=gray!8,
        pattern=north east lines,
        pattern color=gray!65
    ]
        (2.10,2.15)

        % Fronteira superior
        .. controls (2.30,3.30) and (3.50,4.35) .. (5.30,4.75)
        .. controls (6.15,5.05) and (6.95,5.40) .. (7.40,5.60)

        % Borda direita/inferior
        .. controls (7.05,4.80) and (6.70,3.95) .. (6.15,3.20)
        .. controls (5.60,2.65) and (4.95,2.30) .. (4.25,2.10)
        .. controls (3.55,1.90) and (3.10,1.35) .. (3.25,0.80)

        % Retorno até M
        .. controls (2.60,1.15) and (2.25,1.55) .. cycle;


    % Contorno do conjunto factível
    \draw[thick]
        (2.10,2.15)
        .. controls (2.30,3.30) and (3.50,4.35) .. (5.30,4.75)
        .. controls (6.15,5.05) and (6.95,5.40) .. (7.40,5.60)
        .. controls (7.05,4.80) and (6.70,3.95) .. (6.15,3.20)
        .. controls (5.60,2.65) and (4.95,2.30) .. (4.25,2.10)
        .. controls (3.55,1.90) and (3.10,1.35) .. (3.25,0.80)
        .. controls (2.60,1.15) and (2.25,1.55) .. cycle;


    % =========================================================
    % FRONTEIRA EFICIENTE
    % =========================================================

    \draw[
        very thick
    ]
        (2.10,2.15)
        .. controls (2.30,3.30) and (3.50,4.35) .. (5.30,4.75)
        .. controls (6.15,5.05) and (6.95,5.40) .. (7.40,5.60);



    % =========================================================
    % CARTEIRA DE VARIÂNCIA MÍNIMA
    % =========================================================

    \fill (2.10,2.15) circle (2.2pt);

    \node[
        left
    ] at (1.98,2.15)
        {$M$};

    \draw[->]
        (1.4,1.5) -- (1.98,2.05);

    \node[
        left,
        align=right
    ] at (1.45,1.45)
        {Variância\\mínima};


    % =========================================================
    % ALGUMAS CARTEIRAS POSSÍVEIS
    % =========================================================

    \fill (3.15,2.75) circle (1.8pt);
    \node[below right] at (3.15,2.75) {$P_1$};

    \fill (3.75,3.55) circle (1.8pt);
    \node[below right] at (3.75,3.55) {$P_2$};

    \fill (4.65,3.20) circle (1.8pt);
    \node[below right] at (4.65,3.20) {$P_3$};

    \fill (5.55,3.90) circle (1.8pt);
    \node[below right] at (5.55,3.90) {$P_4$};


    % =========================================================
    % PONTO ÓTIMO DO INVESTIDOR
    % =========================================================

    \fill (5.30,4.75) circle (2.3pt);

    \node[
        above left
    ] at (5.30,4.75)
        {$P^*$};


    % =========================================================
    % CURVA DE INDIFERENÇA
    %
    % Representa uma combinação constante de utilidade para
    % determinado grau de aversão ao risco.
    % =========================================================

    \draw[
        thick,
        dashed,
        domain=3.85:6.75,
        samples=80,
        smooth,
        variable=\x
    ]
    plot
    (
        {\x},
        {4.75
         + 0.33*(\x-5.30)
         + 0.28*(\x-5.30)^2}
    );

    \node[
        align=left
    ] at (5.5,6)
        {Investidor 2};


    % =========================================================
    % EXPLICAÇÃO DA REGIÃO FACTÍVEL
    % =========================================================

    \draw[->]
        (7.75,3.15) -- (6.20,3.55);

    \node[
        right,
        align=left
    ] at (7.70,3.10)
        {Carteiras\\possíveis};


    % =========================================================
    % CARTEIRA DE MAIOR RISCO/RETORNO REPRESENTADA
    % =========================================================

    \fill (7.40,5.60) circle (2pt);

    \node[
        above right
    ] at (7.40,5.60)
        {$W$};

    \end{tikzpicture}
\end{comment}

\begin{tikzpicture}[
    x=1.15cm,
    y=1.05cm,
    >=Stealth,
    every node/.style={font=\small}
]

% =========================================================
% PONTOS-CHAVE
% =========================================================

\coordinate (M)     at (2.10,2.15);
\coordinate (Ione)  at (2.55,3.45);
\coordinate (Pstar) at (5.20,4.72);
\coordinate (W)     at (7.30,5.55);

% =========================================================
% EIXOS
% =========================================================

\draw[->, thick]
    (0,0) -- (8.3,0)
    node[right] {Desvio-padrão $(\sigma_p)$};

\draw[->, thick]
    (0,0) -- (0,6.5)
    node[above] {$E(R_p)$};

\node[left, align=center] at (-0.15,5.6)
    {Retorno\\esperado};

% =========================================================
% CONJUNTO DE CARTEIRAS POSSÍVEIS
%
% Ajuste principal:
% - abriu mais a região logo após M;
% - suavizou a concavidade inferior;
% - eliminou o aspecto de "rabo" triangular.
% =========================================================

\path[
    fill=gray!8,
    pattern=north east lines,
    pattern color=gray!65
]
    (M)

    % -------- fronteira eficiente --------
    .. controls (2.10,2.70) and (2.25,3.15) .. (Ione)
    .. controls (2.95,3.85) and (4.45,4.45) .. (Pstar)
    .. controls (6.00,4.95) and (6.75,5.35) .. (W)

    % -------- lateral direita / borda inferior --------
    .. controls (6.95,4.95) and (6.70,4.25) .. (6.35,3.60)
    .. controls (5.90,3.10) and (5.20,2.72) .. (4.45,2.55)
    .. controls (3.95,2.45) and (3.55,2.35) .. (3.35,2.05)
    .. controls (3.18,1.88) and (3.25,1.68) .. (3.42,1.50)
    .. controls (3.00,1.78) and (2.42,1.95) .. (M)
    -- cycle;

% Contorno do conjunto factível
\draw[thick]
    (M)
    .. controls (2.10,2.70) and (2.25,3.15) .. (Ione)
    .. controls (2.95,3.85) and (4.45,4.45) .. (Pstar)
    .. controls (6.00,4.95) and (6.75,5.35) .. (W)
    .. controls (6.95,4.95) and (6.70,4.25) .. (6.35,3.60)
    .. controls (5.90,3.10) and (5.20,2.72) .. (4.45,2.55)
    .. controls (3.95,2.45) and (3.55,2.35) .. (3.35,2.05)
    .. controls (3.18,1.88) and (3.25,1.68) .. (3.42,1.50)
    .. controls (3.00,1.78) and (2.42,1.95) .. (M);

% =========================================================
% FRONTEIRA EFICIENTE
% =========================================================

\draw[very thick]
    (M)
    .. controls (2.10,2.70) and (2.25,3.15) .. (Ione)
    .. controls (2.95,3.85) and (4.45,4.45) .. (Pstar)
    .. controls (6.00,4.95) and (6.75,5.35) .. (W);

% =========================================================
% CARTEIRA DE VARIÂNCIA MÍNIMA
% =========================================================

\fill (M) circle (2.2pt);
\node[left] at (1.97,2.15) {$M$};

\draw[->]
    (1.35,1.45) -- (2.00,2.05);

\node[left, align=right] at (1.40,1.40)
    {Variância\\mínima};

% =========================================================
% CARTEIRAS INTERNAS
% =========================================================

\fill (3.15,2.75) circle (1.8pt);
\node[below right] at (3.15,2.75) {$P_1$};

\fill (3.75,3.55) circle (1.8pt);
\node[below right] at (3.75,3.55) {$P_2$};

\fill (4.60,3.20) circle (1.8pt);
\node[below right] at (4.60,3.20) {$P_3$};

\fill (5.45,3.90) circle (1.8pt);
\node[below right] at (5.45,3.90) {$P_4$};

% =========================================================
% PONTO ÓTIMO P*
% =========================================================

\fill (Pstar) circle (2.3pt);
\node[above left] at (5.15,4.76) {$P^*$};

% =========================================================
% INVESTIDOR 1
%
% Agora é um trecho mais curto e mais suave.
% A ideia é sugerir tangência em Ione sem virar uma haste.
% =========================================================

\draw[thick]
    (1.30,0.80)
    .. controls (1.90,1.45) and (2.15,2.75) .. (Ione)
    .. controls (2.78,4.05) and (2.88,4.95) .. (2.82,5.60);

\fill (Ione) circle (2.2pt);

\node[align=center] at (3.10,5.70) {Investidor 1};

\draw
    (2.92,5.48) -- (2.88,5.00);

% =========================================================
% INVESTIDOR 2
%
% Agora é uma curva curta, tangente visualmente em P*,
% sem parecer uma reta colada na fronteira.
% =========================================================

\draw[thick]
    (3.15,4.18)
    .. controls (4.05,4.35) and (4.70,4.50) .. (Pstar)
    .. controls (5.85,4.95) and (6.55,5.55) .. (6.98,6.12);

\node[align=center] at (5.95,6.18) {Investidor 2};

% =========================================================
% EXPLICAÇÃO DA REGIÃO FACTÍVEL
% =========================================================

\draw[->]
    (7.70,3.20) -- (6.10,3.58);

\node[right, align=left] at (7.66,3.15)
    {Carteiras\\possíveis};

% =========================================================
% CARTEIRA W
% =========================================================

\fill (W) circle (2pt);
\node[above right] at (W) {$W$};

\end{tikzpicture}